\section{Results}
\label{sec:results}

\subsection{Measurement-Chain Audits}
Static position repeatability is 0.04--0.06~mm per axis and rotation
repeatability is 0.008--0.035$^\circ$. At 0.25-px corner noise, Monte Carlo
median position spread is 0.80~mm on validation and 0.72~mm on test; rotation
spread is 0.23$^\circ$ and 0.19$^\circ$, respectively
(Fig.~\ref{fig:uncertainty}). Combining static and corner terms by
(\ref{eq:uncertainty}) gives a representative position
per-axis $U_{\rm op}$ ($k_{\rm c}=2$) of approximately 1.4--1.6~mm for a frame, excluding the
unavailable board/intrinsic tolerances.

\begin{figure}[tb]
\centering
\includegraphics[width=\columnwidth]{reference_uncertainty.pdf}
\caption{Reference characterization. Tilt errors produce millimeter-scale
translation changes while remaining weakly visible in reprojection.}
\label{fig:uncertainty}
\end{figure}

Dynamic holdout depends strongly on frame rate: median position residual is
0.27--0.31~mm at 29~fps, 1.06--1.89~mm at 25~fps, and 3.06--5.19~mm at
10~fps. At the latter rate, median rotation residual is 0.96--1.31$^\circ$.
The tilt sensitivity is 3.4--3.7~mm/deg but only 0.13--0.18~px/deg in
reprojection, so the 1.5-px reprojection gate does not bound pose error.

The rigidity recording yields held-out axis correlations 0.78/0.92/1.00,
magnitude correlation 0.95, and scale 0.99; two earlier recordings made
before the mount was secured fail the test and are excluded. Constraining
PnP rotation to the gyroscope still shifts the test reference by 8.1~mm
median (21.4~mm p90). Training on those contaminated labels raises apparent
0.2-s test $R^2$ from 0.04 to 0.24, consistent with label leakage.

\paragraph{Twin audit.}
All 95 selected inertial sequences and 25 rendered sequences pass the
automated audit. Ideal/noisy recomputation agrees below
$5\times10^{-5}$~g and $10^{-4}$~deg/s per axis. Rendered corner error is
0.252~px median and 0.842~px p90. Fig.~\ref{fig:domain} shows why an initial
twin failed: its median angular rate was three times the real value
(14.1 against 4.3$^\circ$/s) and its translational acceleration was too small.
Retuning these distributions changed single-model real-validation $R^2$ from
0.326 to 0.403.

\subsection{Observable Horizon and the Empirical Error Floor}
\label{sec:limit}
The displacement horizon is the dominant observability parameter
(Fig.~\ref{fig:limit}(a)). For the earlier body-frame Conv--BiGRU (one seed),
the held-out test does not beat zero motion at 0.2~s, and validation $R^2$ rises through 0.177, 0.203, 0.265, and 0.343 for 0.5, 1, 2,
and 3~s, then saturates at 0.339 for 4~s. The selected setting is therefore
$h=3$~s with 2~s context on each side.

At that horizon conventional model changes no longer give a robust gain.
Table~\ref{tab:ablation} lists the fixed-validation choices. Excluding the
deliberately time-reversed data, the Group-A means span 31.89--33.24~mm;
the five-seed reference itself spans 29.80--31.40~mm after the curated
extra-train data are added. Counting the earlier Conv--BiGRU sweeps, more than
forty variants---capacity, regularization, loss, supervision density,
augmentation, pooling, context length, coarse branches, pretraining, input
axes and extra data---show no consistent improvement beyond this seed and
data-selection variation. Time reversal reaches 34.20~mm against 32.56~mm at
the same weight decay, suggesting that the real sensor chain is not
time-symmetric.

\begin{table}[t]
\centering
\caption{PhysNet ablations on the fixed validation split, in millimetres.
Group A screens the architecture on the six trajectory-training recordings
(two seeds; width 96, weight decay 0.1, dense supervision, pre-integration,
lever arm, 1$^\circ$ extrinsic perturbation, 2-s context). Its final row adds
all six rigidity recordings before their mount quality is considered. Group B
uses the final curated pool, which retains only the four rigidity recordings
that pass the held-out rigidity test, and reports the seed count and spread.
The last two rows are five-model ensembles. $^\dagger$Run with weight decay
0.01; compare with the 32.56-mm row.}
\label{tab:ablation}
\scriptsize
\setlength{\tabcolsep}{2.5pt}
\begin{tabular}{lr@{\hspace{3pt}}lr}
\toprule
\multicolumn{4}{l}{\textbf{A. Capacity, loss, augmentation, pretraining, data}} \\
\midrule
Reference & 31.97 & Weight decay 0.01 & 32.56 \\
Width 64 & 32.19 & Dropout 0.2$^\dagger$ & 33.00 \\
End-point loss only$^\dagger$ & 32.83 & No pre-integration ch.$^\dagger$ & 32.55 \\
No lever arm$^\dagger$ & 32.81 & Extrinsic pert. 0 / 3$^\circ$$^\dagger$ & 32.42 / 32.37 \\
Context 3 / 4~s$^\dagger$ & 32.86 / 33.11 & Coarse branch 21 / 31~s & 32.58 / 32.12 \\
Time warp 1.4$\times$ & 31.98 & Amplitude 0.4--1.5$\times$ & 32.48 \\
Velocity loss $w{=}1$ / 3 & 33.24 / 32.58 & Time reversal$^\dagger$ & 34.20 \\
Pretrain, legacy twin & 32.38 & Pretrain, corrected twin & 32.61 \\
Session-balanced & 32.38 & $+$6 rigid (unfiltered) & 31.89 \\
\addlinespace[2pt]
\toprule
\multicolumn{2}{l}{\textbf{B. Zero velocity and uncertainty}} & Seeds & Mean / spread \\
\midrule
\multicolumn{2}{l}{Reference} & 5 & 30.68 \ (29.80--31.40) \\
\multicolumn{2}{l}{Stillness loss only} & 2 & 30.96 \ (30.61--31.32) \\
\multicolumn{2}{l}{Two-pass ZUPT re-anchoring} & 2 & 30.55 \ (30.35--30.75) \\
\multicolumn{2}{l}{Structural gate (\ref{eq:gate})} & 5 & \textbf{30.42} \ (29.86--30.85) \\
\multicolumn{2}{l}{\quad still weight 0.3 / 3} & 2 & 30.23 / 31.12 \\
\multicolumn{2}{l}{\quad $+$ velocity-stream loss} & 2 & 30.42 \ (30.36--30.49) \\
\multicolumn{2}{l}{\quad $+$ two-pass ZUPT} & 2 & 30.92 \ (30.67--31.17) \\
\multicolumn{2}{l}{\quad $+$ heteroscedastic NLL} & 2 & 30.09 \ (29.87--30.31) \\
\multicolumn{2}{l}{Reference, 5-model ensemble} & 5 & 29.91 \\
\multicolumn{2}{l}{Structural gate, 5-model ensemble} & 5 & \textbf{29.67} \\
\bottomrule
\end{tabular}
\end{table}

The plateau has a clear signature (Fig.~\ref{fig:limit}(b)). For the fixed-split
PhysNet ensemble (31.46~mm on validation), the residual correlates with the
reference displacement at $\varrho=-0.88/-0.78/-0.92$ per axis, the regression
slope of prediction on reference is only 0.37/0.48/0.20, and the median
predicted magnitude is 0.44 of the reference. The estimator returns a shrunk,
conditional-mean displacement, which is what (\ref{eq:decomp}) predicts when a
component of the target is not determined by the inputs.

Attributing a quantified share of the residual to ${\bf v}(t_a)$ requires an
independent velocity measurement, which the 10-fps reference cannot supply
(Fig.~\ref{fig:limit}(c)). An affine correction by the mean reference velocity
over the 7-s window removes 14.1\% of the error in-sample but only 4.3\% when
fitted on the other validation recording, and this proxy contains the target
displacement itself. A proxy from the two context segments alone removes
1.8\% in-sample and adds 6.5\% across recordings, and ${\bf v}(t_a)$
differentiated from the reference over $\pm0.3$~s is dominated by PnP noise
(mean $\|{\bf v}(t_a)\Delta t_{ab}\|=75.8$~mm against a mean displacement of
35.6~mm). The real data therefore do not measure how much of the 30-mm residual
the initial velocity causes. The other candidates are residual bias and
gravity leakage, which an ideal IMU in the twin shows to be minor
(Section~\ref{sec:protocol}), reference noise at 10~fps and the anisotropic
training motion. Longer context did not help either: 3-s and 4-s context were
worse (Table~\ref{tab:ablation}), and 756~s of reference motion did not
establish a long-range motion prior.

The shrinkage is not an artefact of the training loss: in post-freeze
fixed-split experiments (two seeds), a near-quadratic loss ($\beta=5$~cm), the
Gaussian NLL as main term and a pause-interpolated gyroscope bias left the
per-axis slopes at 0.42--0.48, 0.47--0.56 and 0.19--0.29 (29.5--31.0~mm), and
scaling the predictions up raised error and ATE (gated: 126 to 152~mm at
$\alpha=1.5$). Averaging gated and non-gated models lowered the CV error from
24.95 to 24.63~mm but not the held-out error below IMUNet (43.89 versus
42.93~mm), as expected of a conditional-mean estimator with missing information.

\begin{figure}[tb]
\centering
\includegraphics[width=\columnwidth]{observability.pdf}
\caption{Observability of the 3-s protocol on validation. The horizon saturates
at 3~s (a). Predictions shrink toward zero; the dashed line is the identity
(b). A velocity proxy removes error only when it includes the target interval
(c).}
\label{fig:limit}
\end{figure}

Session-level fusion did not reduce this error. One sparse least-squares
problem over all node positions with displacement, neighbour-velocity and
zero-velocity edges leaves the validation error at 29.7--29.9~mm against
29.67~mm for independent windows, and the ATE within 2~mm, for every weight
combination tried (variance weighting: below 1.5~mm). The tens of thousands of
displacement edges dominate the few hundred neighbour constraints, so absolute
velocity must be supplied by the acquisition, for example by pauses.

\subsection{Stationary-Boundary Experiment and Acquisition Requirement}
\label{sec:pause}
An epoch correctly identified as stationary supplies the boundary condition
${\bf v}={\bf0}$. Detecting it is reliable: on validation the stillness head
recovers 94.4--100\% of the frames that the reference confirms slower than
3~mm/s, with 0--2.3\% false alarms on frames faster than 8~mm/s.

Fig.~\ref{fig:zupt}(a) measures what such a boundary condition makes available
without a learned model. Between two detected pauses the velocity is pinned at
both ends; given an attitude trajectory, the specific force determines the
displacement, and the two boundary conditions identify one constant
board-frame acceleration offset per interval (2.9--6.1~mg, dominated by the
specific-force magnitude error along gravity). The only interval below 2~s
gives 0.7~mm with reference attitude and 1.2~mm with gyroscope attitude
($n=1$, illustrative only). Pooled 2--4-s errors rise to 53.9 and 104.1~mm, and
longer intervals produce hundreds to thousands of millimetres. A constant tilt
is largely absorbed by the fitted offset, so the divergence indicates
time-varying attitude error, nonconstant bias and noise: a tilt growing with a
gyroscope bias $b_g$ gives $g_0 b_g t^3/6$, whereas a constant one-degree tilt
alone would give $g_0\,\delta\theta\,t^2/2=0.77$~m over 3~s. Part of this growth
comes from the session-level gyroscope bias used throughout: the mean of all
samples below 3~$^\circ$/s also averages slow rotation and removes
0.04--0.19~$^\circ$/s of false bias, whereas strictly still samples give at
most 0.017~$^\circ$/s. With the raw rate the pooled 2--4-s and 8--30-s errors
fall from 104.1~mm and 4.7~m to 78.1~mm and 2.0~m; retrained on the raw rate,
the learned model is unchanged (29.44 against 29.43~mm).

\begin{figure}[tb]
\centering
\includegraphics[width=\columnwidth]{zupt_uncertainty.pdf}
\caption{(a) Zero-velocity-aided integration between detected pauses, pooled
over the training, validation and rigidity recordings. The sub-2-s point is one
interval. The horizontal line is the 43.6-mm held-out error of the gated 3-s
window. (b) Uncertainty coverage on validation, before and after one scale
fitted on the same data.}
\label{fig:zupt}
\end{figure}

The benefit that actually materializes is bounded by how often the instrument
stops: 0.5--2.3 still segments per minute (training 2.3, validation 1.38,
rigidity 1.75, held-out test 0.52), and only 16.5\% of validation windows
contain any still frame. Consistently, using the detection as the structural
update (\ref{eq:gate}) improves the mean single-model validation error by
0.26~mm, inside the 1.6~mm seed range of the reference configuration,
although it narrows that range to 1.0~mm and improves every local
held-out-test statistic while worsening trajectory error
(Table~\ref{tab:ablation} and Section~\ref{sec:estimators}). Supervising the
stillness classifier without using it in the kinematics does not help, and the
two-pass re-anchoring does not combine with the gate. Below 10~mm/s for at
least 0.3~s, 65\% of the validation windows have no stop in
$[t_a-2\,\mathrm{s},t_b]$ (36.5 against 50.8~mm for zero motion), 8\% one
only in the context and 27\% one inside the interval (13.7 against 21.7~mm).

\subsection{Acquisition Protocol in the Twin}
\label{sec:protocol}
To estimate what the frozen protocol's 1--2-s pause every 10--15~s and denser
spacings deliver, synthetic-only models (two seeds) were trained on five
handheld corpora that differ only in the mean stop interval (60/10/20
sequences of 40~s). On this non-periodic
motion PhysNet removes 15--38\% of the zero-motion error
(Table~\ref{tab:protocol}), the range seen on real data, mainly on windows with
a stop inside the interval. An ideal IMU without noise, bias or quantization
(35.1 and 46.8~mm at 4~s and without stops), an 8-s look-back (36.9~mm at
4~s) and a 10- or 30-fps label rate (43.8 and 42.7~mm at 12.5~s) change the
error by at most 2~mm.

\begin{table}[tb]
\centering
\caption{Synthetic-test error (mm) on handheld motion with 1--2-s stops (zero
motion 55.4--57.4~mm). In win.: windows with a stop in $(t_a,t_b]$; PhysNet
with session-level bias removal or the raw gyroscope rate (Raw); +sw.:
stop-anchored switching.}
\label{tab:protocol}
\footnotesize
\setlength{\tabcolsep}{3.2pt}
\begin{tabular}{@{}lcccccc@{}}
\toprule
Interval & In win. & PhysNet & Raw & Raw+sw. & IMUNet & IMUNet+sw. \\
\midrule
4~s    & 72\% & 35.3 & 34.7 & \textbf{30.0} & 57.5 & 33.3 \\
8~s    & 50\% & 42.2 & 39.7 & \textbf{36.7} & 55.7 & 43.2 \\
12.5~s & 43\% & 43.2 & 42.5 & \textbf{40.3} & 55.3 & 46.5 \\
20~s   & 31\% & 46.5 & 44.5 & \textbf{41.8} & 54.2 & 47.8 \\
None   & 15\% & 48.6 & 44.2 & \textbf{44.1} & 55.5 & 55.3 \\
\bottomrule
\end{tabular}
\end{table}

The learned models leave most of the information in a stop unused:
stop-anchored strapdown integration (gravity from the still samples, velocity
drift removed linearly between them) reaches 4.4~mm with the ideal IMU and the
true stops on windows with a stop inside the interval. With the measured noise
and stops detected from the IMU alone ($\|\boldsymbol\omega\|<1^\circ$/s,
$|\|{\bf f}\|-g_0|<0.01g_0$ for 0.5~s) it is accurate only over short spans,
so a window is switched to it only if no stretch longer than $S$ lacks a
bracketing stop ($S$ selected on synthetic validation: 2~s for PhysNet, 3~s for
IMUNet). Switching gives 30.0~mm at 4~s and lets IMUNet, which alone stays
within 2~mm of zero motion, reach 33.3~mm; on real validation it affects 8.5\%
of the windows, almost all static (29.63 to 29.52~mm). With the same estimator
the twin predicts 9\% lower error than pause-free motion (44.1~mm) at the
frozen 10--15-s spacing and 32\% for one firm stop every 4~s.

\subsection{Estimators at the Error Floor}
\label{sec:estimators}
Table~\ref{tab:summary} and Fig.~\ref{fig:benchmark} compare every method
under one protocol. On synthetic data, where the sensor model is exact and the
motion and sensor distributions are controlled by the twin, PhysNet
reduces the error of the best public architecture, IMUNet, by 24\%
(16.69 versus 21.89~mm, $R^2$ 0.864 versus 0.812) and by 8\% on the
out-of-distribution split, with one seventh of the parameters (0.42~M versus
2.83~M); all networks run in 0.13--0.16~ms per window, so latency does not
distinguish them. The synthetic advantage is concentrated in periodic motion:
PhysNet removes 79--87\% of the zero-motion error on the circle, depth, lateral
and mixed families but only 41\% on random-spline motion ($-16$\% on its OOD
sequence), close to the 30\% on real data, so periodic synthetic accuracy
overstates what a window resolves during handheld motion.
Models trained only on the legacy twin fail on real data (PhysNet degrades
least, 90.2~mm). With the gravity-corrected, replay-augmented twin PhysNet
transfers directly (39.6 versus 42.2~mm for zero motion, $R^2$ up to 0.17),
and with a realistic twin (handheld, random-spline and 170 replayed segments)
IMUNet does as well (37.9--38.6~mm; PhysNet 37.5--38.8~mm). Fine-tuning from
either twin, also with motion-class-balanced sampling, brings no gain
(three-seed ensembles: PhysNet 29.29--29.56 against 29.43~mm, IMUNet 31.41
against 30.63~mm).

% Generated by tools/build_summary_table.py; do not edit by hand.
\begin{table*}[t]
\centering
\caption{All methods under one six-axis, 3-D, 3-s displacement protocol. Synthetic columns use synthetic-only models; ``real direct'' evaluates them on the held-out real recording without any real update. ``Val'': fixed-split validation; ``CV'': leave-one-session-out errors of the two-seed fold ensembles with extra-train; ``Held-out test'': CV fold ensembles (16 models) unless marked, with ``fast'' the error in the fastest third; ATE: pose-graph translation error with one anchor per connected component (31 components). Bold marks the best learned value per column; on the held-out test none of the pairwise differences between IMUNet, PhysNet and the blends is significant (Fig.~\ref{fig:benchmark}(b)). Latency is per window for a batch of 16 on one GPU.}
\label{tab:summary}
\footnotesize
\setlength{\tabcolsep}{3.2pt}
\begin{tabular}{lrrrrrrrrrrrr}
\toprule
 & \multicolumn{2}{c}{Synth.\ test} & OOD & Real & Val & CV & \multicolumn{3}{c}{Held-out test} & ATE & Par. & Lat. \\
Method & mm & $R^2$ & mm & direct & mm & mm & mm & $R^2$ & fast & mm & M & ms \\
\midrule
Zero motion & 58.46 & $-0.001$ & 67.31 & 49.7 & 42.23 & 30.42 & 49.67 & $-0.016$ & 90.27 & 119.2 & -- & -- \\
Ridge features & 56.29 & 0.074 & 72.42 & -- & 33.98 & -- & 50.16$^{\dagger}$ & 0.040 & 79.34 & 99.2 & -- & -- \\
Conv--BiGRU, body frame & 26.16 & 0.740 & 43.72 & 90.8 & 30.58 & -- & 46.27$^{\dagger}$ & 0.148 & 77.86 & 102.4 & 0.34 & -- \\
RoNIN-LSTM (adapted) & 58.46 & $-0.001$ & -- & 49.7 & -- & -- & -- & -- & -- & -- & 0.21 & 0.15 \\
RoNIN-ResNet (adapted) & 25.04 & 0.769 & 42.36 & 128.2 & 33.43 & -- & 44.74$^{\dagger}$ & 0.210 & 74.14 & 108.7 & 3.85 & 0.13 \\
TLIO-ResNet (adapted) & 27.40 & 0.726 & 43.44 & 108.2 & 33.90 & -- & 46.15$^{\dagger}$ & 0.185 & 75.53 & -- & 3.85 & 0.14 \\
\midrule
IMUNet (adapted) & 21.89 & 0.812 & 37.20 & 118.8 & 31.90 & 25.77 & 42.93 & 0.250 & 74.20 & 82.2 & 2.83 & 0.14 \\
PhysNet (proposed) & \textbf{16.69} & \textbf{0.864} & \textbf{34.27} & 90.2 & 31.46 & 25.17 & 44.39 & 0.215 & 75.15 & 87.8 & 0.42 & 0.16 \\
PhysNet + stillness gate & -- & -- & -- & -- & \textbf{29.67}$^{\S}$ & 24.95 & 43.55 & 0.252 & 72.67 & 109.8 & -- & -- \\
\midrule
Equal blend, no gate & -- & -- & -- & -- & -- & 24.65 & 42.60 & 0.256 & 74.29 & 81.1 & -- & -- \\
Calibrated blend, no gate & -- & -- & -- & -- & -- & 24.62 & 42.68 & 0.257 & 73.80 & \textbf{80.7} & -- & -- \\
Equal blend, gate & -- & -- & -- & -- & -- & 24.58 & \textbf{42.21} & 0.276 & 72.97 & 90.8 & -- & -- \\
Calibrated blend, gate & -- & -- & -- & -- & -- & \textbf{24.51} & 42.29 & \textbf{0.279} & \textbf{72.10} & -- & -- & -- \\
\bottomrule
\multicolumn{13}{p{0.97\textwidth}}{\footnotesize $^{\dagger}$Single fixed-split model (fine-tuned from the twin where applicable). $^{\S}$Five-seed ensemble trained with the rigidity recordings as extra data; 29.91~mm without the gate on the same data.}
\end{tabular}
\end{table*}


\begin{figure}[tb]
\centering
\includegraphics[width=\columnwidth]{trajectory_results.pdf}
\caption{Pose-graph translation on the held-out recording. The reference
orientation isolates translation. PhysNet and the blend use the non-gated
trajectory operating point.}
\label{fig:trajectories}
\end{figure}


\begin{figure*}[t]
\centering
\includegraphics[width=\textwidth]{estimators.pdf}
\caption{Estimators under one protocol. Architectures are retrained here;
values are not taken from the pedestrian benchmarks. (a) Error reduction
relative to zero motion; a dot marks a combination that was not evaluated,
and the scale is clipped at $-60\%$. (b) Paired differences on the held-out test with
95\% block-bootstrap intervals. Filled markers exclude zero: every learned
model beats zero motion, and none differs significantly from IMUNet.
(c) Leave-one-session-out error, sorted by the zero-motion error.}
\label{fig:benchmark}
\end{figure*}

On real data the two best families do not differ significantly. Gated PhysNet
has the lowest single-family error on fixed validation (29.67~mm; 29.91~mm
without the gate on the same data) and in cross-validation (fold ensembles
24.95 versus 25.77~mm for IMUNet; six of eight sessions won,
Fig.~\ref{fig:benchmark}(c)). On the held-out test the PhysNet$-$IMUNet error
difference is 0.62~mm with a 95\% interval of $[-1.8,\,+3.2]$~mm; averaging the
two families gives the lowest mean error (equal blend, 42.21~mm) and the
calibrated blend ($w=0.65$, $\alpha=1.05$) the highest $R^2$ (0.279), but its
improvement over IMUNet, $[-1.2,\,+2.5]$~mm, also contains zero; with about 38
independent 3-s blocks the recording cannot resolve differences below roughly
3~mm. IMUNet, gated PhysNet and the blends beat zero motion by 6.1--7.5~mm with
intervals that exclude zero (Fig.~\ref{fig:benchmark}(b)); Ridge, the
Conv--BiGRU and the fixed-split RoNIN and TLIO models gain at most 4.9~mm.
PhysNet is thus distinguished by efficiency, the lowest synthetic and OOD
errors and direct transfer at real-data accuracy that matches the larger
models. Learned displacement helps only on the four large-motion recordings
($R^2$ 0.17--0.62), not on the three near-static ones.
Data curation had effects of the same order as architecture: adding the
rigidity recordings moved the IMUNet CV error from 27.17 to 26.38~mm, and
excluding the two that fail the mount test moved the PhysNet validation error
from 31.89 to 30.68~mm. Removing, post freeze, the two trajectory recordings
with the poorest camera--IMU rotation consistency (11.6 and 7.8$^\circ$) moved
the CV error on the remaining six from 27.35 to 26.72~mm and the held-out error
from 43.55 to 43.38~mm, inside its interval.

\subsection{Estimator Performance, Uncertainty, and Physical Consistency}
The heteroscedastic head, evaluated on the two-seed fixed-split ensemble
trained with the NLL term (the selected CV models use $w_\eta=0$), costs no
accuracy (Table~\ref{tab:ablation}). Its raw predicted-plus-ensemble variance
is under-dispersed (32.5\% and 54.5\% coverage at one and two standard
deviations) because it is fitted on the dense stream. The scale $\kappa=4.56$
fitted on validation gives 81.4\% and 94.5\% (Fig.~\ref{fig:zupt}(b)), 81.0\%
and 94.2\% across validation recordings, but only 63.2\% and 85.2\% on the
larger-motion test; stratified, magnitude-scaled and split-conformal
recalibrations did not remove this shortfall. The most and least confident
deciles have 6.9 and 29.7~mm error on validation and 1.7 and 51.5~mm on test.

The sixteen non-gated fold models estimate the camera--IMU lever arm as
$(-22.6\pm3.8,\,-5.2\pm1.9,\,-29.7\pm7.0)$~mm in camera axes (mean $\pm$
standard deviation), a 38-mm offset, whereas the gated models give
$(-20.9,\,-4.1,\,-17.8)$~mm, and removing the term costs only 0.25~mm (32.81
versus 32.56~mm). The estimate is variant-dependent and weakly constrained; it
must be checked against a mechanical measurement, not used as a calibration.

\paragraph{Trajectory error and the operating point.}
The last column of Table~\ref{tab:summary} and Fig.~\ref{fig:trajectories}
report pose-graph ATE on the held-out test. Chessboard dropouts split the
713 reference nodes into 31 connected components; each is anchored to the
reference position of its first node, so the ATE measures drift over
components of 3.7~s on average rather than over the whole 115-s recording,
and zero motion is scored with the same anchors (119.2~mm). The CV fold
ensemble reduces the ATE of the fixed-split IMUNet from 95.2 to 82.2~mm.
PhysNet gives 87.8~mm with one edge per window but 81.1~mm when every
reference frame in its displacement stream replaces the end-point edges
(29\,885 edges); the non-gated blends give 80.7~mm (calibrated) and 81.1~mm
(equal weights).

The stillness gate separates the two metrics: it improves every local
statistic in Table~\ref{tab:summary} while degrading ATE from 87.8 to 109.8~mm.
Both ensembles are equally shrunk (median predicted magnitude 15.8 and 16.4~mm
against 42~mm), so this is not a simple shrinkage effect, and we did not
isolate the mechanism. We report the gated ensemble as the displacement
pseudo-measurement and the non-gated one as the trajectory estimator.

\paragraph{Initially anchored relative pose.}
Replacing the reference orientation in the trajectory solve by bias-corrected
gyroscope integration gives a 4.51$^\circ$ orientation RMSE over the 115-s
held-out recording and a 3-s relative-rotation RMSE of 3.48$^\circ$. With the
non-gated calibrated blend, the relative-pose stream has 80.59-mm translation
ATE and 50.92-mm 3-s RPE, against 119.2~mm and about 60~mm for zero motion.
Attitude is taken from the reference only at the first node and positions use
the same 31 anchors, so this is a relative pose stream with periodic position
anchoring, not a learned or absolute 6-DoF pose.


